\documentclass[11pt]{amsart}
\usepackage[margin=1.1in]{geometry}
\usepackage{amssymb,amsmath,hyperref}
\newtheorem{theorem}{Theorem}
\newtheorem{lemma}[theorem]{Lemma}
\newtheorem{proposition}[theorem]{Proposition}
\theoremstyle{remark}
\newtheorem{remark}[theorem]{Remark}
\newcommand{\C}{\mathbb C}
\newcommand{\R}{\mathbb R}
\newcommand{\CP}{\mathbb{CP}}
\DeclareMathOperator{\tr}{tr}
\DeclareMathOperator{\Herm}{Herm}
\DeclareMathOperator{\diam}{diam}

\title{The complex projector set of $\CP^2$ has Borsuk number nine}
\author{[AUTHOR NAME]}
\date{October 2026}

\begin{document}
\begin{abstract}
OpenAI's family-156 preprint shows that the projector image of $\mathbb{RP}^3$ in
$\R^9$ cannot be covered by ten sets of smaller diameter. The complex analogue,
the projector image $X$ of $\CP^2$ in the eight-dimensional space of trace-one
Hermitian $3\times3$ matrices, is a natural candidate for an eight-dimensional
counterexample. We record that it is not one: $b(X)=9$. The lower bound is the
verbatim complex transfer of the preprint's odd-map argument. The upper bound
is an explicit cover by nine caps, certified on all of $\CP^2$ by a
branch-and-bound computation in outward-rounded interval arithmetic and,
independently, in exact integer arithmetic. The result is a negative one. Since $\CP^n$ for
$n\ge3$ lives in dimension $\ge15$, it closes the route of complex projective
projector sets for counterexamples in dimension $<9$. The mathematics is elementary.
\end{abstract}
\maketitle

\section{Statement}
For a unit vector $u\in\C^3$ write $P_u=uu^*$. Let
\[
 X=\{P_u:\ u\in\C^3,\ |u|=1\}\subset\{A\in\Herm_3(\C):\tr A=1\}\cong\R^8,
\]
with the Frobenius metric $\|A\|_F^2=\tr(A^*A)$. Since $\tr(P_uP_v)=|\langle u,v\rangle|^2$,
\begin{equation}\label{eq:dist}
 \|P_u-P_v\|_F^2=2-2|\langle u,v\rangle|^2 ,
\end{equation}
so $\diam X=\sqrt2$, attained exactly at orthogonal pairs. For a bounded set $Y$ of positive diameter, $b(Y)$ is the least number of subsets of
diameter $<\diam Y$ that cover $Y$.

\begin{theorem}\label{thm:main}
$b(X)=9$. Explicitly, there are nine points $p_1,\dots,p_9\in\C^3$ (listed in
the accompanying file \texttt{P9.txt}) such that the nine closed caps
\[
 C_i=\Big\{P_u\in X:\ |\langle u,p_i\rangle|^2\ \ge\ \tfrac{21}{40}\,|p_i|^2\Big\}
\]
cover $X$ and each has diameter at most $\sqrt{399/200}<\sqrt2$. Conversely
$X$ cannot be covered by eight sets of diameter $<\sqrt2$.
\end{theorem}

Since $X$ spans an $8$-dimensional affine space, $9=8+1$ is exactly Borsuk's
bound, so $X$ is not a Borsuk counterexample.

\section{Lower bound: $b(X)\ge 9$}
This is the complex form of the preprint's strict-gap
reduction (\texttt{geometry.tex}, \texttt{lem:cover-to-map}), of its odd
extension (\texttt{extension.tex}, \texttt{lem:positive-extension},
\texttt{lem:signed-extension}) and of the first sentence of its
Proposition \texttt{thm:label-bound}. The proofs carry over line by line. We
indicate the only changes.

\begin{proposition}\label{prop:lower}
If $X$ is covered by $m$ sets of diameter $<\sqrt2$, then $m\ge 9$.
\end{proposition}
\begin{proof}
\emph{Step 1 (reduction).} As in the preprint (proof of \texttt{lem:cover-to-map}),
enlarge the sets to relatively open sets $U_i$ of diameter $<\sqrt2$ and take a
continuous partition of unity on $\CP^2$ subordinate to their preimages. This gives
continuous $f=(f_1,\dots,f_m):\CP^2\to\Delta^{m-1}$ such that, with
$S(x)=\{i:f_i(x)>0\}$, orthogonal lines $x\perp y$ have $S(x)\cap S(y)=\varnothing$.
(If $i\in S(x)\cap S(y)$ then $P_x,P_y\in U_i$ would be at distance $\sqrt2$
by \eqref{eq:dist}.) Only continuity is needed below.

\emph{Step 2 (extension).} Put $H(v)=|v|^2f([v])$ for $v\ne0$, $H(0)=0$; it is
continuous because $\|H(v)\|_1=|v|^2$. For positive semidefinite $Q\in\Herm_3$ let
\[
 E(Q)=3\int_{S(\C^3)}H(Q^{1/2}u)\,d\mu(u),
\]
with $\mu$ the $U(3)$-invariant probability measure on the unit sphere of $\C^3$.
The map $E$ is continuous, since $Q\mapsto Q^{1/2}$ is continuous on the PSD
cone. It is positively homogeneous and coordinatewise nonnegative. Since
$\int uu^*\,d\mu=\tfrac13I$, it satisfies $\sum_iE(Q)_i=3\int\langle Qu,u\rangle
d\mu=\tr Q$. As in the preprint, $\{i:E(Q)_i>0\}=\bigcup_{x\subseteq \operatorname{ran}Q}S(x)$.
This uses only that $Q^{1/2}$ maps the sphere onto a set spanning every line of
$\operatorname{ran}Q$, and that nonempty open subsets of the sphere have positive measure.

\emph{Step 3 (odd map).} For $A\in\Herm_3$ write $A=A_+-A_-$ with $A_\pm\succeq0$
supported on the orthogonal positive and negative spectral subspaces, and put
$F(A)=E(A_+)-E(A_-)$. Every line in $\operatorname{ran}A_+$ is orthogonal to
every line in $\operatorname{ran}A_-$. By Step~1 the two vectors therefore have disjoint
supports, so $\|F(A)\|_1=\tr A_++\tr A_-=\tr|A|$. Moreover $F$ is continuous and
odd ($A_\pm$ are exchanged under $A\mapsto-A$). Hence $A\mapsto F(A)$ restricts to a
continuous odd map from $\{A\in\Herm_3:\tr|A|=1\}\cong S^{8}$ (as
$\dim_\R\Herm_3=9$) to $\{y\in\R^m:\|y\|_1=1\}\cong S^{m-1}$. By the
Borsuk--Ulam theorem, $m-1\ge8$.
\end{proof}

\begin{remark}[What is and is not analogous]
For real $k\times k$ symmetric matrices the same argument gives $m\ge k(k+1)/2$,
which is again Borsuk's bound $d+1$. The preprint's counterexample comes from
excluding the \emph{equality case} $m=N_4=10$ on $\mathbb{RP}^3$, by a long
degree-theoretic and combinatorial argument (\texttt{topology.tex},
\texttt{six-labels.tex}, \texttt{ten-labels.tex}). Note that $\CP^2$ is the complex
analogue of $\mathbb{RP}^2$ ($k=3$, six labels), where the preprint does not exclude
equality. The complex analogue of the $k=4$ witness is $\CP^3\subset\R^{15}$, which
is irrelevant for the dimension record. The point of Theorem~\ref{thm:main} is only that the
single complex projector set of dimension below $9$ that is not a round sphere,
namely $X$ in $\R^8$, is not a counterexample. ($\CP^1$, $\mathbb{HP}^1$ and $\mathbb{OP}^1$ give round
spheres $S^2,S^4,S^8$ in $\R^3,\R^5,\R^9$, for which $b=d+1$. $\mathbb{HP}^2$ lives in $\R^{14}$.)
\end{remark}

\section{Upper bound: nine caps}
\begin{lemma}[Caps]\label{lem:cap}
Let $p,x,y\in\C^3$ be unit vectors with $|\langle x,p\rangle|^2\ge T$ and
$|\langle y,p\rangle|^2\ge T$, where $T>\tfrac12$. Then
$|\langle x,y\rangle|\ge 2T-1$, and hence
$\|P_x-P_y\|_F^2\le 2-2(2T-1)^2$.
\end{lemma}
\begin{proof}
Write $x=ap+x'$ and $y=bp+y'$ with $x',y'\perp p$. Then $|a|^2\ge T$ and $|x'|^2=1-|a|^2\le1-T$,
and likewise for $y$. Hence $|\langle x,y\rangle|=|a\bar b+\langle x',y'\rangle|\ge|a||b|-|x'||y'|\ge
T-(1-T)$. Apply \eqref{eq:dist}.
\end{proof}

With $T=21/40$ this gives $\diam C_i\le\sqrt{2-2/400}=\sqrt{399/200}\approx1.41244$.
It remains to show that the caps cover $X$. Because the condition is
homogeneous, this is the statement
\begin{equation}\label{eq:cover}
 \forall x\in\C^3\setminus\{0\}:\qquad
 \max_{1\le i\le 9}\ h_i(x)>0,\qquad
 h_i(x)=|\langle x,p_i\rangle|^2-T\,|p_i|^2|x|^2 .
\end{equation}

\begin{proposition}[computer-assisted]\label{prop:cert}
\eqref{eq:cover} holds for the nine points of \texttt{P9.txt} with $T=21/40$,
and also with $T=263/500=0.526$.
\end{proposition}

\emph{Domain.} Every point of $\CP^2$ has a representative with $x_j=1$ and
$|x_k|\le1$ for $k\ne j$: divide by the coordinate of largest modulus. Then
$(\operatorname{Re}x_k,\operatorname{Im}x_k)_{k\ne j}\in[-1,1]^4$. So it suffices to prove
\eqref{eq:cover} on three boxes $[-1,1]^4$ (one per chart $j$). Each
$h_i$ is a real quadratic polynomial in the four free coordinates.

\emph{Primary certificate} (\texttt{cert\_interval.py}). Branch and bound with
16-way splitting. On each box, $\operatorname{Re}$ and $\operatorname{Im}$ of
$\langle x,p_i\rangle$ are linear forms, enclosed in interval arithmetic; their
squares and $|x|^2$ are enclosed termwise. A box is accepted when the lower
endpoint of the enclosure of some $h_i$ is $>0$. Every IEEE-754 operation is
followed by an outward \texttt{nextafter} step, so all enclosures are rigorous.
The data $p_i$ are exact binary64 numbers, and $T|p_i|^2$ is computed exactly as a
rational and then enclosed. For $T=21/40$: $1\,897\,683$ accepted leaves
(minimum half-width $2^{-12}$), under a minute on one core.

\emph{Independent certificate} (\texttt{independent/cert\_exact.py},
standard library only). We use exact integer arithmetic throughout. For a quadratic $H$ and a
box $c\pm r$, it uses the exact expansion $H(c+d)=H(c)+\nabla H(c)\cdot d+d^{\mathsf T}Qd$
and the lower bound
\[
 H(c)-\sum_k|\partial_kH(c)|r_k+\sum_k\min(Q_{kk},0)r_k^2-\sum_{k\ne l}|Q_{kl}|r_kr_l .
\]
Boxes are bisected along the widest coordinate. For $T=21/40$: $16\,876$ leaves,
maximal depth $36$, about 4\,s.

\emph{A symmetric configuration is not enough.} For the Hesse SIC (orbit of
$(0,1,-1)/\sqrt2$ under the Heisenberg group) every overlap of the basis vector $e_1$ with the nine SIC lines is $0$ or $\tfrac12$,
so its maximal overlap is exactly $\tfrac12$ and its caps only reach the borderline $T=\tfrac12$, i.e.\ diameter $\sqrt2$.
A perturbed, computer-found configuration is therefore used.

\emph{Sharpness and control.} Numerical minimization gives
$\min_x\max_i |\langle \hat x,\hat p_i\rangle|^2\approx0.526055$ for these points.
Both certificates succeed at $T=0.526$. At $T=0.53$ the primary code stops with an
explicit violating point
$x=(1,\,0.67578125+0.30859375i,\,-0.26171875+0.08984375i)$, whose normalized
maximal overlap is $0.52950<0.53$. Deleting the ninth point also makes the code
fail immediately, as Proposition~\ref{prop:lower} predicts for any eight caps with $T>\tfrac12$.

\begin{proof}[Proof of Theorem~\ref{thm:main}]
Proposition~\ref{prop:lower} gives $b(X)\ge9$. Proposition~\ref{prop:cert} and
Lemma~\ref{lem:cap} give nine closed caps that cover $X$, each of diameter
$\le\sqrt{399/200}<\sqrt2$.
\end{proof}

\section{Provenance, scope and what is not claimed}
The nine points were found by a heuristic max--min search (soft-min objective,
L-BFGS-B). The search is not part of the proof. The certificate concerns only the
fixed exact numbers in \texttt{P9.txt}.

\begin{itemize}
\item We do not determine the best covering threshold
$\max_{p_1..p_9}\min_x\max_i|\langle\hat x,\hat p_i\rangle|^2$ (it is at least $0.526$),
nor whether symmetric configurations (SIC/MUB/Clifford-equivariant) work. An earlier
exploratory computation suggested that Clifford-equivariant nine-covers have
margin at most $0.005$. That computation is not used here.
\item Nothing about dimensions $\le 8$ in general follows. The note shows only that this
particular candidate set does not give a Borsuk counterexample in $\R^8$.
\item The lower bound is a direct transfer of the preprint's argument. The upper
bound is an elementary covering-radius computation. We do not claim new
methods.
\item There is no Lean formalization.
\end{itemize}

\begin{thebibliography}{9}
\bibitem{OAI156} OpenAI, \emph{A nine-dimensional counterexample to Borsuk's covering assertion},
OpenAI Math Release preprint, family 156, September 23, 2026,
\url{https://github.com/openai/math} (commit \texttt{fd4aeeb2}).
\bibitem{Kalai} G. Kalai, \emph{Some old and new problems in combinatorial geometry I: Around Borsuk's problem},
in: Surveys in Combinatorics 2015, arXiv:1505.04952.
\bibitem{Matousek} J. Matou\v{s}ek, \emph{Using the Borsuk--Ulam theorem}, Springer, 2003.
\bibitem{Hatcher} A. Hatcher, \emph{Algebraic Topology}, Cambridge University Press, 2002.
\bibitem{Moore} R. E. Moore, R. B. Kearfott, M. J. Cloud, \emph{Introduction to Interval Analysis}, SIAM, 2009.
\end{thebibliography}

\bigskip\noindent\small Developed with AI assistance (Anthropic Claude, Opus 5.5) under
[AUTHOR NAME]'s direction; independent adversarial reviews were run by separate AI
sessions and are not human peer review.
\end{document}
